\documentclass{jsarticle}
\usepackage[dvipdfmx]{graphicx}
\usepackage{amsmath}
\usepackage{amssymb}
\begin{document}
\title{カルーツァ・クライン空間とシュワルツシルト半径 \\
そのまわりのブラックホールのエントロピー値 \\
量子群と重力場方程式における多様体積分}
\author{Masaaki Yamaguchi and My son in acasic record of \\
this assented world on heavens gate.}
\date{}
\maketitle

    カルーツァ・クライン空間の方程式は、
    $$ ds^2 = g_{\mu\nu}(x)dx^{\mu}dx^{\nu} + \psi^2(x)(dx^2 + \kappa^2
    A_{\mu}(x)dx^{\nu})^2 $$
    この式は、
    $$ ds^2 = -N(r)^2dt^2 + \psi^2(x)(dr^2 + r^2d\theta^2) $$
    $$ ds^2 = -dt^2 + r^{-8\pi Gm}(dr^2 + r^2d\theta^2) $$
    とまとまり、
    $$ dx^2 = g_{\mu\nu}(x)(g_{\mu\nu}(x)dx^2 - dxg_{\mu\nu}(x)) $$
    $$ dx = (g_{\mu\nu}(x)^2dx^2 - g_{\mu\nu}(x)dx g_{\mu\nu}(x))^{1 \over 2} $$
    と反重力と正規部分群の経路和となり、　
    $$ \pi(\chi, x) = i \pi (\chi, x)f(x) - f(x)  \pi (\chi, x) $$
    と基本群にまとまる。
    シュワルツシルト半径は、
    $$ ds^2 = -(1 - {r_s \over r})c^2 dt^2 + {1 \over {{1 - {r_s \over r}}}} 
    dr^2 + r^2 d \theta^2 + r^2 \sin^2 \theta d \psi^2 $$
    これは、カルーツァ・クライン空間と同型となる。
    一般相対性理論は、
    $$ R^{\mu\nu} + {1 \over 2}g_{ij}\Lambda = \kappa T^{\mu\nu} $$
    この多様体積分は、
    $$ \int \kappa T^{\mu\nu}\mathrm{dvol} = \int  \left(R^{\mu\nu} + 
    {1 \over 2}g_{ij}\Lambda\right)\mathrm{dvol}  $$
    ガンマ関数のおける大域的微分多様体は、これと同型により、
    $$ \int \Gamma \cdot {d \over {d \gamma}}{\Gamma dx_m} \le \int \Gamma
    dx_m + {d \over d \gamma}\Gamma $$
    $$ = \int \Gamma(\gamma)^{'}dx_m $$
    オイラーの定数の多様体積分は、
    $$ \int \left({{\int {1 \over x^s}}dx - \log x} \right)\mathrm{dvol} $$
    この解は、シュワルツシルト半径と同型より、
    $$ = e^f - e^{-f} \le {e^f + e^{-f}} $$
    次元の単位は多様体より、大域的微分多様体とオイラーの定数
    の多様体積分の加群分解は、オイラーの公式の三角関数の虚数の度解より、
    $$ {d \over df}F + \int C dx_m = 2 (\cos (i x \log x) 
    - i \sin (i x \log x)) $$
    すべては、大域的微分多様体の重力と反重力方程式に行き着き、
    $$ {d \over df}F \ge {d \over df}\int \int{1 \over (x \log x)^2}dx_m
    + {d \over df}\int \int{1 \over (y \log y)^{1 \over 2}}dy_m $$
    と求まる。

    一般相対性理論における多様体積分とオイラーの定数の多様体積分
    が同型と言えることを上の式たちは述べている。

    このJones多項式が、金融市場の相場価格を決めるオプション方程式であり、
    流体力学による熱エントロピー値の熱エネルギー量が、
    いろんな機材に使われている。
    量子コンピュータにおける論理素子の回路の設計にも使われている。
    人工知能の論理素子と因数分解における量子暗号にも使われている。

    この一般相対性理論における多様体積分がJones多項式となり、
    金融市場で使われることを読んだのが、イスラエル国家らしい。

    $$ H\Psi = \bigoplus {\left(i \hbar^{\nabla} \right)}^{\oplus L} $$
    $$ = {1 \over 2}i e^{i \Hat{H}} $$
    $$ (i \hbar)^{'} = (- e^{i \Hat{H}})^{'} $$
    $$ = - i e^{i \Hat{H}} $$
    $$ \psi(x) = e^{-i \Hat{H}t}, \bigoplus (i \hbar ^{\nabla})^{\oplus L} 
    = {{{1 \over 2} e^{i \Hat{H}}}^{(-i e^{i \Hat{H}})}} $$
    $$ = ({1 \over 2}f)^{-i f}$$
    $$ = ({1 \over 2})^{-i f} \cdot e^{-x \log x} \cdot (f)^{i} $$
    $$ = \int e^{- x} x^{t - 1}dx, {d \over d \gamma}\Gamma = e^{-x \log x}  $$
     
    $$ f = x, i = t, {1 \over 2} = a, \bigoplus a^{-t x}x^t [I_m]
     \cong \int e^{-x} x^{t - 1}dx $$
     $$ |\psi(t)\rangle_s = e^{-i\Hat{H} t}|\Psi\rangle_{H},
    \Hat{A}_s = \Hat{A}_H(0) $$
    $$ |\Psi(t)\rangle_s \to {d \over dt} $$
    $$ i {d \over dt}|\psi (t) \rangle_s = \Hat{H}|\psi(t)\rangle_s $$
    $$ \langle\Hat{A}(t)\rangle = \langle\Psi(t)|\Hat{A}(0)|\Psi(t)\rangle  $$
    $$ {d \over dt}\Hat{A} = {1 \over i}[\Hat{A},H] $$
    $$ \Hat{A}(t) = e^{i\Hat{H}t}\Hat{A}(0)e^{-i\Hat{H}t} $$
    $$ \lim_{\theta \to 0} \begin{pmatrix} \sin \theta \\
    \cos \theta \end{pmatrix} \begin{pmatrix} \theta & 1 \\
    1 & \theta \end{pmatrix} \begin{pmatrix} \cos \theta \\
    \sin \theta \end{pmatrix} = \begin{pmatrix} 1 & 0 \\
    0 & - 1 \end{pmatrix} $$
    $$ f^{-1}(x) x f(x) = I^{'}_m, I^{'}_m = [1,0] \times [0,1] $$

    $$ {x + y} \ge \sqrt{x y} $$
    $$ {x^{{1 \over 2} + iy} \over e^{x \log x}} = 1 $$
    $$ \mathcal{O}(x) = \nabla_i \nabla_j \int e^{{2 \over m}\sin \theta 
    \cos \theta } \times {{N \mathrm{mod}(e^{x \log x})} \over {O(x)(x +
    \Delta |f|^2)^{1 \over 2}}} $$
    $$ x \Gamma(x) = 2 \int|\sin 2 \theta|^2 d \theta, \mathcal{O}(x)
    = m(x)[D^2 \psi] $$
    $$ i^2 = (0,1) \cdot (0,1), |a||b|\cos \theta = - 1 $$
    $$ E = \mathrm{div}(E, E_1) $$
    $$ \left({\{f,g\}} \over [f,g] \right) = i^2, E = mc^2, I^{'} = i^2 $$

    この式たちは、微分幾何の量子化から計算方法としての形式の
    大域的微分多様体の大域的部分積分が,大域的偏微分方程式としての式の
    形からわかるのも、微分幾何の量子化の仕組みとしての計算式から
    導けられるのも、すべては未来の私の子の情報から読み取れたと
    言える。
 
    夏目漱石が、JRが出来ていない明治時代の蒸気機関車も出来ていないときに、
    電車が広島駅と名古屋駅について書いていて、これは、三四郎であり、
    蒸気機関車は草枕であり、両者とも九州を始めとして、
    男を強調して、このあとに、イギリスに渡り、アカシックレコードを
    ほのめかして、統合失調症の精神分裂病と双極性障害ともども、
    助けることを書いている人とわかったときに、すきー、と言えます。

    弘法大師様は、自身が御影になり、御影石になり、この世の、
    見られているということを、自身が神隠しの役割をして、
    消しているというが、犯罪が消えていないし、手に入って離れるも
    ないし、でも、念力を般若心経に込めて、代々伝えていることを
    見ると、自身の役割をしっかりしている、日本にとっての宝と言える。

    アインシュタイン博士は、最近になって知ったことであるが、
    子どものために、つけていて、読まれないことを実践して、
    役割が去っていったので、去るべきときと言っていた。
    歳晩年に看護師さんと関係を持ったらしいのも教えてもらえて、
    ファンを大切にするいい人とおもった。


    ルパン三世の能力は、私の若かりし頃に姉２人が隠し持っていた、
    漫画本と小説の場所が手に取るようにわかる能力と同じということです。
    空気みたいにわかっていて、なぜその場所にあるかがわかるのが不思議
    な能力だったことです。高等学校になって、私から消えた能力です。

    
\end{document}
